\documentclass[a4paper,12pt]{article}
\usepackage[T2A]{fontenc}
\usepackage[utf8]{inputenc}
\usepackage[russian]{babel}
\usepackage{amsmath,amssymb,mathtools}
\renewcommand{\familydefault}{\sfdefault}
\usepackage{icomma}
\usepackage[a4paper,margin=20mm,top=18mm,bottom=20mm]{geometry}
\usepackage{microtype}
\usepackage{tikz}
\usetikzlibrary{calc,arrows.meta,positioning,shapes.geometric}
\usepackage{tabularx,booktabs,longtable}
\usepackage{enumitem}
\usepackage{parskip}
\usepackage{fancyhdr}
\usepackage{setspace}
\usepackage{xcolor}
\usepackage{hyperref}

\definecolor{accent}{HTML}{2F6F6D}
\definecolor{darkaccent}{HTML}{214F4D}
\definecolor{textgray}{HTML}{2F3335}
\definecolor{softgray}{HTML}{70777A}
\definecolor{rulegray}{HTML}{D6DCDA}

\hypersetup{hidelinks}
\setlength{\parindent}{0pt}
\setlength{\parskip}{5pt}
\onehalfspacing
\color{textgray}
\setlist[itemize]{leftmargin=6mm,itemsep=2pt,topsep=3pt}
\setlist[enumerate]{leftmargin=8mm,itemsep=5pt,topsep=4pt}

\pagestyle{fancy}
\fancyhf{}
\renewcommand{\headrulewidth}{0pt}
\fancyfoot[C]{\small\color{softgray}04.10.26 \quad\textbullet\quad Марк Гукин \quad\textbullet\quad \thepage}

\newcommand{\sectionline}{\par\vspace{2pt}\noindent\color{rulegray}\rule{\linewidth}{0.5pt}\color{textgray}\par\vspace{5pt}}
\newcommand{\formula}[1]{\[\displaystyle #1\]}
\newcommand{\markrule}[1]{\textcolor{accent}{\textbf{#1}}}
\newcommand{\exhead}[1]{\par\medskip\textcolor{darkaccent}{\textbf{#1}}\par}

\begin{document}

% ---------- TITLE ----------
\thispagestyle{empty}
\begin{center}
\vspace*{20mm}
{\color{softgray}\large МАТЕМАТИКА}\par
\vspace{12mm}
{\fontsize{29}{34}\selectfont\bfseries\color{darkaccent} Чек-лист}\par
\vspace{5mm}
{\fontsize{20}{25}\selectfont\bfseries Степени и корни: преобразование выражений}\par
\vspace{12mm}
{\large Практический алгоритм, свойства, типичные ошибки и тренировка}\par
\vfill
{\large 04.10.26}\par
\vspace{7mm}
{\normalsize Лёвин Артём Александрович}\par
{\normalsize\bfseries эксклюзивно для Марка Гукина}\par
\vspace{12mm}
\end{center}

\begin{tikzpicture}[remember picture,overlay]
  \draw[accent!55,line width=1.2pt]
    ([xshift=18mm,yshift=16mm]current page.south west)
      .. controls +(42mm,12mm) and +(-48mm,-7mm) ..
    ([xshift=-18mm,yshift=20mm]current page.south east);
\end{tikzpicture}

\clearpage

% ---------- CORE CHECKLIST ----------
{\LARGE\bfseries\color{darkaccent}1. Что нужно уметь видеть сразу}\par
\sectionline

Главная идея занятия: \markrule{сначала привести выражение к удобному степенному виду}, затем применять свойства степеней. Большинство громоздких примеров упрощаются после того, как основания становятся одинаковыми.

\exhead{1.1. Переводите составные числа к удобному основанию}
\formula{49=7^2,\qquad 27=3^3,\qquad 8=2^3,\qquad 42=6\cdot 7.}
Если рядом уже есть степень числа $7$, то число $49$ выгодно сразу заменить на $7^2$.

\exhead{1.2. Переводите дроби и десятичные числа в степени}
\formula{0{,}5=\frac12=2^{-1},\qquad 0{,}25=\frac14=2^{-2}.}
Это особенно полезно, когда в выражении уже присутствуют степени двойки.

\exhead{1.3. Переводите корни в дробные показатели}
Для положительного основания:
\formula{\sqrt[n]{a}=a^{1/n},\qquad \sqrt[n]{a^m}=a^{m/n}.}
Например,
\formula{\sqrt[3]{2}=2^{1/3},\qquad \sqrt[12]{2}=2^{1/12}.}
Если корень сам возводится в степень или вложен в другой корень, удобно использовать:
\formula{(\sqrt[n]{a})^k=a^{k/n},\qquad \sqrt[k]{\sqrt[n]{a}}=\sqrt[kn]{a}.}
Условия существования корней в действительных числах сохраняются.

\exhead{1.4. Сначала работайте со структурой, затем считайте}
Не пытайтесь сразу вычислять большие числа. Сначала найдите одинаковые основания, одинаковые показатели, степени степеней и произведения в степенях.

\vspace{5pt}
{\LARGE\bfseries\color{darkaccent}2. Свойства степеней}\par
\sectionline

Для выражений, где все действия определены:
\begin{align*}
a^m\cdot a^n&=a^{m+n},\\
\frac{a^m}{a^n}&=a^{m-n},\qquad a\ne 0,\\
(a^m)^n&=a^{mn},\\
(ab)^n&=a^n b^n,\\
a^{-n}&=\frac1{a^n},\qquad a\ne0,\\
a^0&=1,\qquad a\ne0.
\end{align*}

\markrule{Удобный обратный ход:} если показатели одинаковы, можно объединить основания:
\formula{a^r b^r=(ab)^r.}
Для дробных показателей безопаснее применять это правило к положительным основаниям.

% ---------- EXAMPLES ----------
{\LARGE\bfseries\color{darkaccent}3. Пять опорных примеров}\par
\sectionline

\exhead{Пример 1. Привести к одному основанию}
\formula{7^{4/9}\cdot49^{5/18}}
Так как $49=7^2$,
\begin{align*}
7^{4/9}\cdot49^{5/18}
&=7^{4/9}\cdot(7^2)^{5/18}\\
&=7^{4/9}\cdot7^{5/9}\\
&=7^{9/9}=7.
\end{align*}

\exhead{Пример 2. Корни как степени}
\formula{\left(\frac{\sqrt[3]{2}\,\sqrt[4]{2}}{\sqrt[12]{2}}\right)^2}
Внутри скобок:
\formula{2^{1/3+1/4-1/12}=2^{4/12+3/12-1/12}=2^{1/2}.}
Следовательно,
\formula{\left(2^{1/2}\right)^2=2.}

\exhead{Пример 3. Десятичная дробь в степени}
\formula{\frac{0{,}5^{\sqrt{10}-1}}{2^{-\sqrt{10}}}}
Так как $0{,}5=2^{-1}$,
\begin{align*}
\frac{(2^{-1})^{\sqrt{10}-1}}{2^{-\sqrt{10}}}
&=2^{-(\sqrt{10}-1)+\sqrt{10}}\\
&=2^1=2.
\end{align*}

\exhead{Пример 4. Одинаковый показатель у разных оснований}
\formula{\frac{6^{\sqrt3-1}\,7^{\sqrt3-1}}{42^{\sqrt3-2}}}
Объединяем числитель:
\formula{6^{\sqrt3-1}7^{\sqrt3-1}=42^{\sqrt3-1}.}
Тогда
\formula{42^{(\sqrt3-1)-(\sqrt3-2)}=42.}

\exhead{Пример 5. Выражение с переменной}
При $m\ne0$:
\formula{\frac{7(m^5)^6+11(m^3)^{10}}{(3m^{15})^2}}
После раскрытия степеней:
\formula{\frac{7m^{30}+11m^{30}}{9m^{30}}=\frac{18m^{30}}{9m^{30}}=2.}

% ---------- DOMAIN / ERRORS ----------
{\LARGE\bfseries\color{darkaccent}4. Ограничения и типичные ошибки}\par
\sectionline

\begin{itemize}
  \item \textbf{Нельзя делить на ноль.} В формулах с делением фиксируйте условие $a\ne0$.
  \item \textbf{$a^0=1$ только при $a\ne0$.} Выражение $0^0$ в школьных преобразованиях не используется как определённое число.
  \item \textbf{Отрицательная степень означает обратное число:} $a^{-n}=1/a^n$ при $a\ne0$.
  \item \textbf{Для корня чётной степени} в действительных числах подкоренное выражение должно быть неотрицательным.
  \item \textbf{Для отрицательной дробной степени} основание должно удовлетворять условиям дробной степени и одновременно не быть нулём.
  \item \textbf{Минус перед скобками меняет все знаки:}
  \formula{p-(q-1)=p-q+1.}
  Именно на этом шаге часто теряется единица в показателе.
  \item \textbf{Не складывайте основания:} из $a^m\cdot a^n$ получается $a^{m+n}$, а не $(2a)^{m+n}$.
  \item \textbf{Не складывайте показатели при сложении выражений:} $a^m+a^m=2a^m$, но не $a^{2m}$.
\end{itemize}

\exhead{Мини-проверка перед ответом}
\begin{enumerate}
  \item Все составные числа, дроби и корни переведены в удобный вид?
  \item Скобки со степенью раскрыты корректно?
  \item Одинаковые основания объединены?
  \item При делении показатели вычтены с правильными знаками?
  \item Проверены ограничения: знаменатели, нулевое основание, корни чётной степени?
  \item Если получилась отрицательная степень, приведён ли ответ к обычной форме, если это требуется?
\end{enumerate}

\clearpage

% ---------- FLOWCHART ----------
\thispagestyle{plain}
\begin{center}
{\LARGE\bfseries\color{darkaccent}Алгоритм: как упрощать степенное выражение}\par
\vspace{10mm}
\end{center}

\begin{center}
\begin{tikzpicture}[
  node distance=8mm,
  startstop/.style={ellipse,draw=accent,line width=.9pt,minimum width=42mm,minimum height=10mm,align=center,font=\small},
  process/.style={rectangle,rounded corners=2mm,draw=accent,line width=.9pt,minimum width=112mm,minimum height=12mm,text width=102mm,align=center,font=\small},
  decision/.style={diamond,aspect=3.1,draw=accent,line width=.9pt,text width=45mm,align=center,font=\small,inner sep=1.5mm},
  arrow/.style={-{Stealth[length=2.8mm]},line width=.9pt,draw=darkaccent}
]
\node[startstop] (start) {Получено степенное выражение};
\node[process,below=of start] (scan) {Найдите составные числа, дроби, корни, одинаковые основания и одинаковые показатели};
\node[decision,below=of scan,yshift=-1mm] (rewriteq) {Можно сделать основания удобнее?};
\node[process,below=of rewriteq,yshift=-1mm] (rewrite) {Разложите составные числа; представьте дроби и корни как степени};
\node[process,below=of rewrite] (brackets) {Раскройте степень степени и степень произведения: показатели перемножаются};
\node[process,below=of brackets] (combine) {Объедините одинаковые основания: при умножении показатели складываются, при делении вычитаются};
\node[process,below=of combine] (simplify) {Упростите показатель: аккуратно раскройте скобки, приведите дроби к общему знаменателю, сократите противоположные слагаемые};
\node[process,below=of simplify] (check) {Проверьте ограничения и приведите отрицательные степени к обычной форме, если это требуется};
\node[startstop,below=of check] (finish) {Запишите окончательный ответ};

\draw[arrow] (start)--(scan);
\draw[arrow] (scan)--(rewriteq);
\draw[arrow] (rewriteq)--node[right,font=\small]{да}(rewrite);
\draw[arrow] (rewrite)--(brackets);
\draw[arrow] (brackets)--(combine);
\draw[arrow] (combine)--(simplify);
\draw[arrow] (simplify)--(check);
\draw[arrow] (check)--(finish);
\draw[arrow] (rewriteq.east) -- ++(24mm,0) node[above,font=\small,pos=.45]{нет} |- (brackets.east);
\end{tikzpicture}
\end{center}

\clearpage

% ---------- TRAINING ----------
{\LARGE\bfseries\color{darkaccent}5. Тренировочный блок — Путь А}\par
\sectionline

\textbf{Формат: 10 задач.} Решайте письменно. В каждом задании показывайте ключевые преобразования: замена основания, раскрытие степени степени, сложение или вычитание показателей. В задачах с переменной учитывайте указанные ограничения.

\begin{enumerate}
  \item Упростите:
  \[
  7^{4/9}\cdot49^{5/18}.
  \]

  \item Упростите:
  \[
  \left(\frac{\sqrt[3]{2}\,\sqrt[4]{2}}{\sqrt[12]{2}}\right)^2.
  \]

  \item Упростите:
  \[
  \frac{0{,}5^{\sqrt{10}-1}}{2^{-\sqrt{10}}}.
  \]

  \item Упростите:
  \[
  \frac{6^{\sqrt3-1}\,7^{\sqrt3-1}}{42^{\sqrt3-2}}.
  \]

  \item Вычислите:
  \[
  \frac{3^{\sqrt5+10}\cdot3^{5-\sqrt5}}{3^{10}}.
  \]

  \item Вычислите:
  \[
  \sqrt[3]{\sqrt[4]{2^{24}}}.
  \]

  \item Вычислите:
  \[
  (0{,}25)^{-3/2}.
  \]

  \item При $m\ne0$ упростите:
  \[
  \frac{7(m^5)^6+11(m^3)^{10}}{(3m^{15})^2}.
  \]

  \item При $a\ne0$ упростите:
  \[
  \frac{(2a^4)^3a^{-5}}{8a^7}.
  \]

  \item При $x\ne0$ упростите:
  \[
  \frac{\left(\sqrt[3]{7x^6}\right)^2}{\sqrt[3]{49}\,x^4}.
  \]
\end{enumerate}

\clearpage

% ---------- ANSWERS ----------
{\LARGE\bfseries\color{darkaccent}6. Ответы}\par
\sectionline

\renewcommand{\arraystretch}{1.35}
\begin{center}
\begin{tabularx}{0.72\linewidth}{@{}>{\bfseries}c X@{}}
\toprule
№ & Ответ \\
\midrule
1 & $7$ \\
2 & $2$ \\
3 & $2$ \\
4 & $42$ \\
5 & $243$ \\
6 & $4$ \\
7 & $8$ \\
8 & $2$ \\
9 & $1$ \\
10 & $1$ \\
\bottomrule
\end{tabularx}
\end{center}

\vspace{10mm}
\textbf{Контрольный ориентир.} Если ответ не совпал, ищите ошибку в одном из четырёх мест: перевод к общему основанию, умножение показателей при степени степени, знак при вычитании показателей, арифметика дробей в показателе.

\vfill
\begin{center}
{\color{softgray}\small Лёвин Артём Александрович эксклюзивно для Марка Гукина}
\end{center}

\end{document}